\subsection{从质量反证到最小修复与回归门禁}
\label{sec:repair-method}
对可机械核验的失败，进一步形成“固定源码与输入、保存最小失败见证、只改相关表达式、重跑目标及既有条件”的修复闭环。
修复不是由模型评分变高来验收，而须使原失败输入通过，并检查修复前通过的目标检查是否退化。
原始包字节始终保留；修复副本具有新指纹，不借用原版本回执。
本轮两处补丁由源码分析确定，不是训练出的模型或自动发现所有缺陷的算法。
执行前已查看相关代码并固定补丁，因此定位试验属于已知候选问题的针对性验证，不能写成盲检准确率。

独立验收只覆盖登记子契约。例如画像工具的行数与缺失数正确，不代表DQS所有维度、类型推断或全部建议正确。
IQR任务要求使用调用者给定阈值；其精确有理数参照独立于原实现的浮点百分位函数。
缺失定义也必须条件化：把NA、空格当空值可能符合原工具约定，却不符合“仅空字符串为空”的新任务。
后一种情形是任务不适配，应请求替代实现或报告边界，不能笼统标为工具坏。

网络改进相应区分两种干预：改变路由规则，以及修复/补充候选实现。
若加入修复节点后，普通最短路径与证据路径都达到相同成功数，收益应归于候选能力改善，不能全归于复杂路由。
否则仅过滤失败节点还可能通过增加弃权降低完整求解率。
